\documentclass[10pt, twocolumn, landscape, a4paper]{article}

\usepackage[margin=1cm]{geometry}
\usepackage[utf8]{inputenc}
\usepackage[spanish,es-noshorthands]{babel}
\usepackage{amssymb, amsmath, amsbsy}
\usepackage{tasks}
\usepackage{enumerate}
\usepackage[pdftex]{hyperref}
\hypersetup{pdftitle={Ejemplos con conjuntos},pdfauthor={Luis Carrillo Gutiérrez}}

\fontfamily{cmss}\selectfont
\renewcommand{\familydefault}{\sfdefault}
\pagestyle{empty}

\begin{document}
\paragraph*{Conjuntos}

\begin{itemize}
\item{Indicar el valor de verdad de las siguientes afirmaciones :
\begin{enumerate}[I.]
	\item{Si $n(A) = 2$ y $n(B) = 2$; el número máximo de subconjuntos de $A \cup B$ es 8}
	\item{Si $A \cap B = \varnothing$; entonces : $A = \varnothing \wedge B = \varnothing$}
	\item{Si $A - B = \{3;\,4\}$ y $B - A = \{5;\,6\}$; El mínimo número de elementos de \\ $A \cup B$ es 5}
\end{enumerate}

\begin{tasks}(5)
	\task{FFF}\task{VVF}\task{FFV}\task{VFF}\task{VVV}
\end{tasks}}
\item{Se tiene dos conjuntos comparables, cuyos cardinales se diferencian en $3$. La diferencia de los cardinales de sus conjuntos potencias es $112$. La suma de los cardinales de los dos conjuntos comparables es:
\begin{tasks}(5)\task{15}\task{13}\task{12}\task{11}\task{10}
\end{tasks}}
\item{Hallar el valor de verdad de las siguientes proposiciones : 
\begin{enumerate}[I.]
\item{Si $n(A) = 4$\quad y \quad$n(B) = 5$; entonces el número máximo de elementos de : \\ $C = P(A) \cup P(B)$ es $48$}
\item{Si $M = \{n^{\,2} - 1 / n \in \mathbb{Z},\; 4 \leq n^{\,2} \leq 9\}$; entonces : $n(M) = 2$}
\item{Si $A = \varnothing$;\quad $B = \varnothing$; entonces : $n\Big[P(A \cup B)\Big] = \varnothing$}
\end{enumerate}

\begin{tasks}(5)
\task{VVF}\task{VFV}\task{FFV}\task{FVF}\task{FFF}
\end{tasks}
}
\item{Dados dos conjuntos disjuntos $A$ y $B$ de ``$n$'' y ``$n + 2$'' elementos respectivamente. Se sabe que el número de subconjuntos de uno de ellos es mayor que el otro en $3072$. Hallar el valor de ``$n - 1$''.
\begin{tasks}(5)
\task{10}\task{9}\task{8}\task{7}\task{6}
\end{tasks}}

\item{Indicar el valor de verdad de las siguientes proposiciones :
\begin{enumerate}[I.]
\item{$\{\varnothing\} \cup \{0\} = \{0\}$}
\item{$\varnothing \cap \{\varnothing\} = \varnothing$}
\item{$\varnothing \cup \{0\} = \{0\}$}
\item{$\{2;\,\varnothing\} \in \{t2;\,\varnothing;\,3\}$}
\end{enumerate}
Los respectivos valores de verdad son : 
\begin{tasks}(3)
\task{FFFF}\task{VFVF}\task{FVVF}\task{VVVF}\task{VVVV}
\end{tasks}}
\end{itemize}
\end{document}